\documentclass[11pt]{article}

% Change "review" to "final" to generate the final (sometimes called camera-ready) version.
% Change to "preprint" to generate a non-anonymous version with page numbers.
\usepackage[review]{acl}

% Standard package includes
\usepackage{times}
\usepackage{latexsym}

% For proper rendering and hyphenation of words containing Latin characters (including in bib files)
\usepackage[T1]{fontenc}
% For Vietnamese characters
% \usepackage[T5]{fontenc}
% See https://www.latex-project.org/help/documentation/encguide.pdf for other character sets

% This assumes your files are encoded as UTF8
\usepackage[utf8]{inputenc}

% This is not strictly necessary, and may be commented out,
% but it will improve the layout of the manuscript,
% and will typically save some space.
\usepackage{microtype}

% This is also not strictly necessary, and may be commented out.
% However, it will improve the aesthetics of text in
% the typewriter font.
\usepackage{inconsolata}

%Including images in your LaTeX document requires adding
%additional package(s)
\usepackage{graphicx}

\usepackage{amsmath}
\usepackage{amssymb}
\usepackage{amsfonts}
\usepackage{algorithm}
\usepackage{algpseudocode}
\usepackage{booktabs}       % For professional tables
\usepackage{xcolor}         % For colored text
\usepackage{pifont}         % For checkmark/cross symbols
\usepackage{array}          % For advanced table formatting
\usepackage{subcaption}
\usepackage{dblfloatfix}
\usepackage{placeins}
\usepackage{tabularx}
\usepackage{float}

% Define colored symbols
\newcommand{\cmark}{\textcolor{green!60!black}{\ding{51}}}  % Green checkmark
\newcommand{\xmark}{\textcolor{red!70!black}{\ding{55}}}    % Red cross
\newcommand{\wmark}{\textcolor{orange!80!black}{!}}         % Orange warning (optional)

% If the title and author information does not fit in the area allocated, uncomment the following
%
%\setlength\titlebox{<dim>}
%
% and set <dim> to something 5cm or larger.

\title{Recovering Instruction Following Loss in Quantized LLMs via Attention Head Compensation}

% Author information can be set in various styles:
% For several authors from the same institution:
% \author{Author 1 \and ... \and Author n \\
%         Address line \\ ... \\ Address line}
% if the names do not fit well on one line use
%         Author 1 \\ {\bf Author 2} \\ ... \\ {\bf Author n} \\
% For authors from different institutions:
% \author{Author 1 \\ Address line \\  ... \\ Address line
%         \And  ... \And
%         Author n \\ Address line \\ ... \\ Address line}
% To start a separate ``row'' of authors use \AND, as in
% \author{Author 1 \\ Address line \\  ... \\ Address line
%         \AND
%         Author 2 \\ Address line \\ ... \\ Address line \And
%         Author 3 \\ Address line \\ ... \\ Address line}

\author{First Author \\
  Affiliation / Address line 1 \\
  Affiliation / Address line 2 \\
  Affiliation / Address line 3 \\
  \texttt{email@domain} \\\And
  Second Author \\
  Affiliation / Address line 1 \\
  Affiliation / Address line 2 \\
  Affiliation / Address line 3 \\
  \texttt{email@domain} \\}

%\author{
%  \textbf{First Author\textsuperscript{1}},
%  \textbf{Second Author\textsuperscript{1,2}},
%  \textbf{Third T. Author\textsuperscript{1}},
%  \textbf{Fourth Author\textsuperscript{1}},
%\\
%  \textbf{Fifth Author\textsuperscript{1,2}},
%  \textbf{Sixth Author\textsuperscript{1}},
%  \textbf{Seventh Author\textsuperscript{1}},
%  \textbf{Eighth Author \textsuperscript{1,2,3,4}},
%\\
%  \textbf{Ninth Author\textsuperscript{1}},
%  \textbf{Tenth Author\textsuperscript{1}},
%  \textbf{Eleventh E. Author\textsuperscript{1,2,3,4,5}},
%  \textbf{Twelfth Author\textsuperscript{1}},
%\\
%  \textbf{Thirteenth Author\textsuperscript{3}},
%  \textbf{Fourteenth F. Author\textsuperscript{2,4}},
%  \textbf{Fifteenth Author\textsuperscript{1}},
%  \textbf{Sixteenth Author\textsuperscript{1}},
%\\
%  \textbf{Seventeenth S. Author\textsuperscript{4,5}},
%  \textbf{Eighteenth Author\textsuperscript{3,4}},
%  \textbf{Nineteenth N. Author\textsuperscript{2,5}},
%  \textbf{Twentieth Author\textsuperscript{1}}
%\\
%\\
%  \textsuperscript{1}Affiliation 1,
%  \textsuperscript{2}Affiliation 2,
%  \textsuperscript{3}Affiliation 3,
%  \textsuperscript{4}Affiliation 4,
%  \textsuperscript{5}Affiliation 5
%\\
%  \small{
%    \textbf{Correspondence:} \href{mailto:email@domain}{email@domain}
%  }
%}

\begin{document}
\maketitle
\begin{abstract}
% Model quantization is essential for deploying large language models (LLMs), 
% yet quantized models often fail at instruction-following in unpredictable ways-generating responses in wrong languages, ignoring formatting constraints, or 
% producing degenerate outputs. We present \textbf{Deep Attention Stimulation (DAS)}, a training-free intervention that 
% selectively compensates degraded attention most affected by quantization. Inspired by neural stimulation in cognitive neuroscience, where localized interventions target dysfunctional neural activity, we identify critical attention heads by analyzing activation differences between full-precision and quantized models on failure cases. Through qualitative analysis on 10 carefully selected 
% samples from IFEval, we show that injecting small corrective signals 
% (alpha-scaled activation differences) into these heads can recover 
% instruction-following behavior: a 4-bit GPTQ Qwen2.5-7B model shifts from 
% producing German output to correct English, and eliminates repetitive 
% degeneration. Our pilot study suggests that quantization-induced 
% instruction-following failures are localized to specific attention heads, 
% opening avenues for interpretable and targeted post-quantization correction methods.

Model quantization is essential for deploying large language models (LLMs). However, quantized models often exhibit unpredictable failures in instruction following, including unintended language switching, violation of formatting constraints, and degenerative generation.
We present \textbf{Deep Attention Stimulation (DAS)}, a training-free intervention that selectively compensates attention heads most disrupted by quantization. Inspired by targeted neural stimulation in cognitive neuroscience, DAS identifies critical attention heads by analyzing activation differences between full-precision and quantized models on instruction-following failure cases. Through qualitative analysis on 10 carefully selected samples from IFEval, we show that injecting small corrective signals into these heads can recover instruction-following behavior.
In particular, a 4-bit GPTQ Qwen2.5-7B-Instruct model recovers correct English output and avoids repetitive degeneration under moderate stimulation. Our pilot study suggests that quantization-induced instruction-following failures are localized to a small subset of attention heads rather than uniformly distributed. These findings highlight the potential of interpretable and targeted post-quantization repair mechanisms.
\end{abstract}

\section{Introduction}
\label{sec:intro}

% ---------------- Figure 1: neuroscience analogy (safe placeholder) ----------------
% \begin{figure*}[t]
%   \centering
%   \includegraphics[width=0.8\textwidth]{figures/figure1-1.png}
%   \caption{\textbf{Neuroscience-inspired Deep Attention Stimulation.} 
%   \textbf{Left:} In cognitive neuroscience, localized neural dysfunction 
%   is diagnosed through neural recordings and remediated via targeted 
%   stimulation. 
%   \textbf{Right:} DAS applies an analogous approach to quantized LLMs. 
%   Comparing attention patterns between full-precision (FP16, top blue) 
%   and quantized models (GPTQ, middle red) identifies degraded heads 
%   (e.g., Head $h$ shows reduced attention to instruction tokens). 
%   Applying $\alpha$-scaled corrections ($\alpha \cdot \Delta\mathbf{O}$) 
%   to these heads restores attention patterns (bottom green) and recovers 
%   instruction-following capability. Heatmaps visualize attention weights 
%   where brighter regions indicate higher attention to instruction tokens.}
%   \label{fig:neuro_analogy}
% \end{figure*}
\begin{figure*}[t]
  \centering
  \includegraphics[width=0.8\textwidth]{figures/figure1-1.png}
  \caption{\textbf{Neuroscience-inspired Deep Attention Stimulation.} 
  Analogous to diagnosing localized neural dysfunction through recordings 
  and correcting it via targeted stimulation (left panel), DAS compares 
  attention patterns between full-precision (FP16, top blue) and quantized 
  models (GPTQ, middle red) to identify degraded heads with reduced 
  instruction token attention (right panel). Applying $\alpha \cdot \Delta\mathbf{O}$ 
  corrections to these heads restores attention patterns (bottom green) 
  and recovers instruction-following capability. Heatmaps visualize attention 
  weights where brighter regions indicate higher attention.}
  \label{fig:neuro_analogy}
\end{figure*}

Large language models (LLMs) have achieved remarkable success in instruction following, enabling users to specify complex constraints on language, format, style, and structure through natural language prompts~\citep{ouyang2022training,wei2022finetuned}. 
However, deploying such models in real-world, resource-constrained environments almost inevitably requires aggressive compression, particularly low-bit quantization~\citep{frantar2022gptq,lin2023awq}. 
While quantized LLMs can often preserve perplexity and many downstream task scores, it has been increasingly observed that instruction-following behavior degrades in subtle yet severe ways after quantization~\cite{qin2025empirical}. 
In practice, quantized models may respond in unintended languages, violate strict formatting constraints, or fall into repetitive and degenerate generations-failure modes that are poorly captured by standard evaluation metrics but critically impact usability.

Instruction following is a fragile and multifaceted capability, relying on precise coordination between semantic understanding, constraint tracking, and long-range control~\citep{zhou2023ifeval}. 
Benchmarks such as IFEval explicitly evaluate whether models adhere to lexical, structural, and stylistic constraints expressed in prompts, revealing failures that may be invisible under conventional accuracy-centric evaluations. 
Notably, quantization-induced instruction-following failures can be highly inconsistent: a quantized model may succeed on most prompts while catastrophically failing on a small subset, even when its full-precision counterpart performs correctly. 
This raises a fundamental question: \emph{are instruction-following failures under quantization the result of global degradation across the model, or do they stem from localized disruptions within specific internal components?}

Inspired by diagnostic and intervention techniques in cognitive neuroscience-particularly electroencephalography (EEG) and targeted neural stimulation-we approach this question from a mechanistic perspective~\citep{dayan2001theoretical,buzsaki2006rhythms}. 
In neuroscience, specific cognitive impairments are often traced back to dysfunction in specific brain regions rather than uniform degradation across the whole brain, enabling localized stimulation to improve or even restore functions. 
Drawing an analogy to transformer-based LLMs, we view attention heads as functional “neural clusters” whose coordinated activity underpins instruction-following behaviors. 
We hypothesize that low-bit quantization selectively disrupts a small subset of attention heads that are disproportionately responsible for maintaining instruction adherence, and that compensating these heads may recover lost functionality.

Motivated by this hypothesis, we introduce \textbf{Deep Attention Stimulation (DAS)}, 
a training-free intervention that identifies and compensates attention heads most affected by quantization.
By comparing activation patterns between full-precision and quantized models on prompts where instruction following fails, DAS isolates a small set of critical heads exhibiting the largest deviations.
During inference, we inject small corrective signals derived from these activation differences directly into the affected heads to correct the "functional loss" during quantization.
Through qualitative analysis on selected failure cases from IFEval~\citep{zhou2023ifeval}, we demonstrate that this localized intervention can recover instruction-following behavior in a 4-bit GPTQ Qwen2.5-7B-Instruct model~\citep{wang2023qwen,frantar2022gptq}, correcting erroneous language switching and mitigating degenerative repetition.
\textbf{Our current analysis on 10 samples demonstrates initial findings that instruction-following failures in quantized LLMs are not uniformly distributed but arise from localized disruptions}
(selected as representative failure cases; see Appendix~\ref{app:sample-selection}),
motivating further large-scale investigation into targeted, interpretable post-quantization repair methods.

Figure~\ref{fig:neuro_analogy} summarizes the neuroscience-inspired motivation and the two-stage diagnostic stimulation workflow of DAS.

% % ---------------- Figure 1: neuroscience analogy (safe placeholder) ----------------
% \begin{figure*}[t]
%   \centering
%   \includegraphics[width=0.8\textwidth]{figures/neuro_analogy.png}
%   \caption{\textbf{Neuroscience-inspired Deep Attention Stimulation.} 
%   \textbf{Left:} In cognitive neuroscience, localized neural dysfunction 
%   (marked with red region and !) is diagnosed via EEG recordings and 
%   corrected through targeted electrical stimulation. 
%   \textbf{Right:} DAS applies an analogous approach to quantized LLMs. 
%   Comparing full-precision (FP16) and quantized (GPTQ) attention patterns 
%   identifies degraded heads (e.g., Head $h$ shows diffused attention away 
%   from instruction tokens like ``lowercase''). These heads are compensated 
%   with $\alpha$-scaled corrective signals ($\alpha \cdot \Delta O$), 
%   restoring instruction-following capability. Both systems demonstrate 
%   how localized dysfunction can be repaired through targeted intervention 
%   rather than system-wide modification.}
%   \label{fig:neuro_analogy}
% \end{figure*}

\section{Related Work}
\label{sec:related}

\subsection{Quantization of LLMs}
Quantization has become a common strategy for improving the efficiency of LLMs by reducing numerical precision and thus lowering memory footprint and inference cost. In transformer-based architectures, post-training quantization (PTQ) has been widely adopted due to its training-free nature~\citep{frantar2022gptq,xiao2023smoothquant, lin2023awq, yao2022zeroquant}. In parallel, quantization-aware training (QAT) integrates quantization effects directly into the training process~\citep{dettmers2023qlora, liu2024llm}, improving task performance under low-precision constraints at additional training and memory cost.


Recent work has extended these techniques to LLMs with billions of parameters, addressing challenges such as highly non-uniform weight distributions and activation outliers. Empirical results show that models including LLaMA, Qwen, and GPT-style architectures can often be compressed to 4-bit and in some cases even 3-bit precision while largely preserving perplexity and downstream task accuracy~\citep{dettmers2022llmint8,frantar2022gptq,lin2023awq}. Low-bit quantization can maintain strong aggregate performance across a wide range of tasks, while exhibiting heterogeneous effects across evaluation dimensions, including alignment-related metrics~\citep{jin2024comprehensive, dettmers2023spqr, kharinaev2025investigating}. Nevertheless, most existing evaluations focus on aggregate accuracy or language understanding benchmarks, leaving open questions about how low-bit quantization affects higher-level behaviors such as instruction-following and reasoning.

\subsection{Evaluation on Instruction Following}
Instruction-following has emerged as a central capability of modern large language models, particularly after instruction tuning and alignment pro~\citep{zhang2025iheval, bang2023multitask}. To evaluate this behavior, recent benchmarks move beyond traditional language modeling metrics and assess whether models can correctly interpret, follow, and adhere to explicit natural language instructions. Representative benchmarks such as IFEval~\citep{zhou2023ifeval}, AlpacaEval~\citep{li2023alpacaeval}, and AlignBench~\citep{liu2024alignbench} focus on constraint satisfaction, format compliance, and rule-following behaviors, providing a more behavior-oriented view of model performance compared to standard perplexity or accuracy-based evaluations.


While prior work on quantization primarily evaluates model quality using aggregate metrics such as perplexity or standard language understanding benchmarks, these evaluations do not fully capture the behaviors exhibited by instruction-tuned large language models in real-world settings~\citep{jin2024comprehensive, xia-etal-2023-training}.
\subsection{Attention Heads in Transformers}

A substantial body of recent work has shown that attention heads in transformer models exhibit significant functional diversity, with different heads specializing in distinct linguistic, syntactic, or positional patterns~\citep{voita2019analyzing,clark2019does,rai2024practical}. More broadly, inspired by diagnostic practices in cognitive neuroscience-where brain neural signals (e.g., EEG) are used to identify behaviorally salient events-recent work in mechanistic interpretability has adopted functional and diagnostic perspectives to analyze internal components of neural networks \citep{kriegeskorte2008representational, nanda2023progress}. 





\section{Method}
\label{sec:method}

\begin{figure}[t]
  \centering
  \begin{subfigure}[b]{0.48\columnwidth}
    \centering
    \includegraphics[width=\textwidth]{figures/figure2.png}
    \caption{Diagnostic phase}
    \label{fig:das_diagnostic}
  \end{subfigure}
  \hfill
  \begin{subfigure}[b]{0.5\columnwidth}
    \centering
    \includegraphics[width=\textwidth]{figures/figure2b.png}
    \caption{Intervention phase}
    \label{fig:das_intervention}
  \end{subfigure}
  \caption{\textbf{Overview of Deep Attention Stimulation (DAS).}
  \textbf{(a)} Comparing FP16 and quantized attention patterns identifies 
  degraded heads with reduced instruction token attention.
  \textbf{(b)} Inference applies $\alpha \cdot \Delta\mathbf{O}$ corrections 
  to these heads via forward hooks, restoring capability without retraining.}
  \label{fig:das_overview}
\end{figure}

\subsection{Problem Setup}

Let $\mathcal{M}_{\text{FP}}$ denote a full-precision instruction-tuned language
model and $\mathcal{M}_{\text{Q}}$ its post-training quantized counterpart (e.g.,
GPTQ-4bit).
Given a set of instruction-following prompts $\mathcal{D}$, we focus on
\emph{failure cases} where $\mathcal{M}_{\text{FP}}$ satisfies all instruction
constraints while $\mathcal{M}_{\text{Q}}$ violates at least one.
Our objective is to diagnose which internal attention heads are most affected by
quantization and to apply a training-free, inference-time intervention that
restores instruction-following behavior. Figure~\ref{fig:das_overview} illustrates the two-stage DAS framework.

\subsection{Locating Instruction-Bearing Tokens}

Instruction-following constraints are often explicitly stated in the prompt
(e.g., ``all lowercase'', ``no commas'', ``respond in English'').
For a prompt $x$, we approximate the \emph{instruction-bearing token set}
$\mathcal{I}(x)$ using lightweight heuristics such as keyword matching and
pattern rules.
This step provides a coarse localization of instruction-related tokens and does
not attempt to fully parse or interpret the instruction semantics.

\subsection{Head-Level Instruction Attention}
\label{sec:head_attention}

Consider a transformer layer $\ell \in \{1,\ldots,L\}$ and attention head
$h \in \{1,\ldots,H\}$.
Let $A^{(\ell,h)}_{\text{FP}}(x)$ and $A^{(\ell,h)}_{\text{Q}}(x)$ denote the
attention weight matrices (after softmax) produced by $\mathcal{M}_{\text{FP}}$
and $\mathcal{M}_{\text{Q}}$ on input $x$, respectively.
We quantify how strongly head $(\ell,h)$ attends to instruction-bearing tokens
by aggregating attention mass onto $\mathcal{I}(x)$:
\begin{equation}
a^{(\ell,h)}(x) =
\frac{1}{|\mathcal{T}(x)|}
\sum_{t \in \mathcal{T}(x)}
\sum_{i \in \mathcal{I}(x)}
A^{(\ell,h)}[t,i],
\label{eq:instruction_attention}
\end{equation}
where $\mathcal{T}(x)$ denotes all query token positions within the prompt and
generated sequence for input $x$.
For brevity, we omit the model subscript when the context is clear.

\subsection{Diagnosing Quantization-Induced Deviation}
\label{sec:diagnosis}

Let $\mathcal{F} \subset \mathcal{D}$ denote the set of instruction-following
failure cases.
We define the \emph{instruction-saliency deviation} of head $(\ell,h)$ as:
\begin{equation}
\Delta a^{(\ell,h)} =
\frac{1}{|\mathcal{F}|}
\sum_{x \in \mathcal{F}}
\left(
a^{(\ell,h)}_{\text{FP}}(x)
-
a^{(\ell,h)}_{\text{Q}}(x)
\right).
\label{eq:delta_attention}
\end{equation}

% Positive values indicate increased attention to instruction tokens after
% quantization, while negative values indicate suppression.
Positive values indicate reduced instruction attention under quantization, while negative values indicate amplification.
In this pilot study, we rank heads by the magnitude
$|\Delta a^{(\ell,h)}|$, capturing both amplification and attenuation effects
induced by quantization.
We then select the top-$k$ most affected heads:
\begin{equation}
\mathcal{H}^\star =
\text{Top-}k
\left(
\left\{
|\Delta a^{(\ell,h)}|
\right\}_{\ell,h}
\right).
\label{eq:critical_heads}
\end{equation}

\subsection{Deep Attention Stimulation}
\label{sec:das}

\paragraph{Correction target.}
We intervene at the level of individual attention head outputs, which allows
localized modification of model behavior without altering weights or retraining.

Let $O^{(\ell,h)}_{\text{FP}}(x)$ and $O^{(\ell,h)}_{\text{Q}}(x)$ denote the
context vectors produced by head $(\ell,h)$ at layer $\ell$ under the
full-precision and quantized models, respectively.
We define the per-head corrective signal as:
\begin{equation}
\Delta O^{(\ell,h)}(x)
=
O^{(\ell,h)}_{\text{FP}}(x)
-
O^{(\ell,h)}_{\text{Q}}(x).
\label{eq:delta_output}
\end{equation}

In our implementation, $\Delta O^{(\ell,h)}(x)$ is computed and cached per failure
sample during the diagnostic phase.

\paragraph{Inference-time intervention.}
During inference with the quantized model, DAS injects an $\alpha$-scaled
correction into the affected heads:

\begin{equation}
\begin{aligned}
\widetilde{O}^{(\ell,h)}(x) &= O^{(\ell,h)}_{\text{Q}}(x) + \alpha \cdot \Delta O^{(\ell,h)}(x), \\
&\qquad \forall\, (\ell,h) \in \mathcal{H}^\star.
\end{aligned}
\label{eq:das}
\end{equation}

This intervention is implemented using forward hooks in the attention module and
does not modify model parameters.
Heads not in $\mathcal{H}^\star$ remain unchanged.

\paragraph{Clarification of $\Delta O$ and Controlled Evaluation.}
In our current study, $\Delta O^{(\ell,h)}(x)$ is computed and applied
in a prompt-specific manner.
That is, for each diagnostic failure case $x$,
we cache the attention head output difference
between the full-precision and quantized models.
During inference, the corrective signal is applied
only when evaluating the same prompt $x$.
This design ensures that sequence length and token positions
match exactly between diagnosis and intervention,
allowing a clean before-and-after comparison
that isolates the effect of attention-level stimulation.

While this prompt-specific formulation does not generalize
to unseen inputs, it provides a controlled setting
for analyzing whether quantization-induced failures
can be repaired through targeted attention intervention.
Developing prompt-agnostic or averaged correction schemes
is left for future work.

\subsection{Algorithmic Summary}

Algorithm~\ref{alg:das} summarizes the complete DAS procedure.
The method consists of two stages:
(1) an offline diagnostic phase that identifies critical heads using a small set
of failure cases, and
(2) an online intervention phase that applies cached corrections during
inference with the quantized model.

% \begin{algorithm}[t]
% \caption{Deep Attention Stimulation (DAS)}
% \label{alg:das}
% \begin{algorithmic}[1]
% \Require Full-precision model $\mathcal{M}_{\text{FP}}$, quantized model
% $\mathcal{M}_{\text{Q}}$, prompts $\mathcal{D}$, number of heads $k$, strength
% $\alpha$
% \Ensure Inference with DAS corrections

% \State \textbf{// Offline Diagnostic Phase}
% \State Identify failure set $\mathcal{F} \subset \mathcal{D}$
% \For{each $x \in \mathcal{F}$}
%   \State Locate instruction tokens $\mathcal{I}(x)$
%   \State Extract attention matrices
%   $A^{(\ell,h)}_{\text{FP}}(x)$ and $A^{(\ell,h)}_{\text{Q}}(x)$
%   \State Compute $a^{(\ell,h)}(x)$ via Eq.~\eqref{eq:instruction_attention}
% \EndFor
% \State Compute $\Delta a^{(\ell,h)}$ via Eq.~\eqref{eq:delta_attention}
% \State Select critical heads $\mathcal{H}^\star$ via
% Eq.~\eqref{eq:critical_heads}
% \State Cache $\Delta O^{(\ell,h)}(x)$ for $(\ell,h)\in\mathcal{H}^\star$

% \State \textbf{// Online Intervention Phase}
% \For{each test prompt $x'$}
%   \State Run $\mathcal{M}_{\text{Q}}(x')$ with forward hooks enabled
%   \For{each $(\ell,h) \in \mathcal{H}^\star$}
%     \State Retrieve cached $\Delta O^{(\ell,h)}(x')$
%     \State Apply correction via Eq.~\eqref{eq:das}
%   \EndFor
%   \State Generate output with corrected activations
% \EndFor
% \end{algorithmic}
% \end{algorithm}

\begin{algorithm}[t]
\caption{Deep Attention Stimulation (DAS)}
\label{alg:das}
\begin{algorithmic}[1]
\Require Full-precision model $\mathcal{M}_{\text{FP}}$, quantized model
$\mathcal{M}_{\text{Q}}$, prompts $\mathcal{D}$, number of heads $k$, strength
$\alpha$
\Ensure Inference with DAS corrections (prompt-specific)

\State \textbf{// Offline Diagnostic Phase}
\State Identify failure set $\mathcal{F} \subset \mathcal{D}$ where $\mathcal{M}_{\text{FP}}$ passes and $\mathcal{M}_{\text{Q}}$ fails
\For{each $x \in \mathcal{F}$}
  \State Locate instruction tokens $\mathcal{I}(x)$
  \State Extract attention matrices $A^{(\ell,h)}_{\text{FP}}(x)$ and $A^{(\ell,h)}_{\text{Q}}(x)$
  \State Compute $a^{(\ell,h)}(x)$ via Eq.~\eqref{eq:instruction_attention}
\EndFor
\State Compute $\Delta a^{(\ell,h)}$ via Eq.~\eqref{eq:delta_attention}
\State Select critical heads $\mathcal{H}^\star$ via Eq.~\eqref{eq:critical_heads}
\For{each $x \in \mathcal{F}$}
  \State Cache $\Delta O^{(\ell,h)}(x)$ for all $(\ell,h)\in\mathcal{H}^\star$
\EndFor

\State \textbf{// Online Intervention Phase (controlled evaluation)}
\For{each prompt $x \in \mathcal{F}$} \Comment{current setup evaluates the same prompts as diagnosis}
  \State Run $\mathcal{M}_{\text{Q}}(x)$ with forward hooks enabled
  \For{each $(\ell,h) \in \mathcal{H}^\star$}
    \State Retrieve cached $\Delta O^{(\ell,h)}(x)$
    \State Apply correction via Eq.~\eqref{eq:das}
  \EndFor
  \State Generate output with corrected activations
\EndFor
\end{algorithmic}
\end{algorithm}

\section{Experiments}
\label{sec:experiments}

\subsection{Experimental Setup}

We evaluate Deep Attention Stimulation (DAS) on Qwen2.5-7B-Instruct,
an instruction-tuned large language model,
quantized to 4-bit precision using GPTQ~\citep{frantar2022gptq}.
All experiments use greedy decoding with fixed generation parameters,
including maximum generation length and decoding strategy,
to control for decoding variability and isolate the effect
of attention-level intervention.
This evaluation protocol follows prior analyses of quantized language models
that emphasize controlled decoding for mechanistic comparison
\citep{dettmers2022llmint8,lin2023awq}.

Evaluation prompts are drawn from IFEval~\citep{zhou2023ifeval},
a benchmark designed to test fine-grained instruction-following constraints
such as output language, formatting, casing, and structural requirements.
From the full benchmark, we identify failure cases in which
the full-precision model satisfies all verifiable constraints,
while the GPTQ-4bit model violates at least one.
We select ten representative samples for detailed evaluation,
covering diverse failure modes including unintended language switching,
formatting violations, constraint omission,
and degenerative repetition.

Crucially, this setup enables a controlled comparison between
instruction-following behavior before repair,
as exhibited by the quantized baseline,
and after attention-based repair,
under identical model weights and decoding conditions.

\subsection{Baseline and Intervention Settings}

We evaluate DAS under three stimulation strengths,
$\alpha \in \{0.0, 5.0, 10.0\}$.
When $\alpha = 0.0$, no corrective signal is injected,
and the model corresponds exactly to the unmodified GPTQ-4bit baseline.
We treat this setting as the before-repair condition.
The setting $\alpha = 5.0$ represents moderate stimulation,
intended to compensate attention heads most affected by quantization.
The setting $\alpha = 10.0$ represents excessive stimulation
and is used to examine the effect of over-intervention.

All before and after comparisons use the same quantized model
and identical decoding parameters,
with the only difference being the presence or absence
of attention-based stimulation.
This design ensures that observed differences
can be attributed to the intervention rather than confounding factors.

\subsection{Manual Evaluation Protocol}

Instruction-following failures often involve subtle qualitative violations
that are not reliably captured by automatic metrics 
\cite{lin2022truthfulqa, li2023evaluating}.
For this reason, all evaluations are conducted by manual inspection.
For each prompt, we derive a checklist of verifiable constraints
from the IFEval instruction,
including output language, casing, formatting,
punctuation, and structural requirements.

An output is labeled as recovered
if all constraints are satisfied,
partially recovered if at least one previously violated constraint
is corrected while others remain violated,
and failed otherwise.
Although IFEval provides automatic constraint-level scores,
prior work has shown that automatic metrics
can overlook instruction-level errors and qualitative failure modes
in large language model outputs
\citep{lin2022truthfulqa,li2023evaluating}.
Manual evaluation therefore remains necessary
for accurately assessing instruction adherence in complex prompts.

This protocol supports both controlled quantitative comparison
across samples and detailed qualitative case analysis.

\subsection{Qualitative Before and After Repair Case Studies}

\begin{figure*}[t]
  \centering
  \includegraphics[width=\textwidth]{figures/figure3.png}
  \caption{Representative qualitative example illustrating instruction-following
  behavior before and after repair.
  Under GPTQ-4bit quantization with $\alpha = 0.0$,
  the model violates the instruction by producing German output.
  With moderate stimulation at $\alpha = 5.0$,
  the model recovers correct English output.
  Excessive stimulation at $\alpha = 10.0$
  destabilizes generation and leads to degenerative repetition.}
  \label{fig:qualitative}
\end{figure*}

We first examine representative qualitative cases
to illustrate how DAS alters instruction-following behavior.
Figure~\ref{fig:qualitative} presents a canonical cross-lingual failure case.
Before repair, the quantized model violates the language constraint,
despite satisfying the instruction under full precision.
After moderate stimulation,
the same model satisfies all instruction constraints.
In contrast, excessive stimulation destabilizes generation,
indicating that DAS functions as a repair mechanism
rather than a general performance enhancer.

% \begin{table*}[b!]
% \centering
% \small
% \begin{tabular}{@{}c>{\raggedright}p{3.5cm}p{3.8cm}p{3.8cm}p{3.8cm}@{}}
% \toprule
% \textbf{ID} & \textbf{Instruction Constraint} & \textbf{GPTQ ($\alpha=0$)} & \textbf{DAS ($\alpha=5$)} & \textbf{DAS ($\alpha=10$)} \\
% \midrule
% 57  & Zen-like style + lowercase + 3 bullets & \xmark{} Historical narrative style  & \cmark{} Abstract Zen metaphors      & \xmark{} Unstable symbols \\[0.3em]
% \textbf{97}  & \textbf{English lowercase letter}   & \xmark{} \textbf{German output (scored 1.0!)} & \cmark{} \textbf{Correct English}   & \xmark{} Repetitive text \\[0.3em]
% 74  & Limerick (AABBA) with highlights    & \xmark{} Broken rhyme scheme         & \xmark{} Still broken rhyme          & \xmark{} Collapsed output \\[0.3em]
% 79  & Letter frequency + no forbidden words & \cmark{} All constraints satisfied  & \cmark{} Maintained success          & \xmark{} Off-task meta-text \\[0.3em]
% 206 & 1930s jazz style + keyword ``rate''   & \cmark{} Proper song structure       & \xmark{} Severe repetition loop      & \xmark{} Continued repetition \\
% \bottomrule
% \end{tabular}
% \caption{Representative instruction-following behaviors under different stimulation strengths.
% \textbf{Sample 97} demonstrates cross-lingual recovery (German $\rightarrow$ English) 
% despite automatic metrics assigning perfect score (1.0) to the German output.
% \textbf{Sample 57} shows stylistic constraint recovery.
% \textbf{Sample 206} illustrates that DAS can degrade initially correct outputs.
% \textbf{Sample 74} remains unrecovered across all settings.
% \cmark{} indicates constraint satisfaction; \xmark{} indicates violation.
% Overall: 2/7 baseline failures recovered (28.6\%), 1/3 baseline passes degraded.}
% \label{tab:qualitative}
% \end{table*}

\begin{table*}[b!] 
\centering
\small
\begin{tabularx}{\textwidth}{@{} c >{\raggedright\arraybackslash}X >{\raggedright\arraybackslash}X >{\raggedright\arraybackslash}X >{\raggedright\arraybackslash}X @{}}
\toprule
\textbf{ID} & \textbf{Instruction Constraint} & \textbf{GPTQ ($\alpha=0$)} & \textbf{DAS ($\alpha=5$)} & \textbf{DAS ($\alpha=10$)} \\
\midrule
57  & Zen-like style + lowercase + 3 bullets & \xmark{} Historical narrative style  & \cmark{} Abstract Zen metaphors       & \xmark{} Unstable symbols \\[0.3em]
\textbf{97}  & \textbf{English lowercase letter}   & \xmark{} \textbf{German output (scored 1.0!)} & \cmark{} \textbf{Correct English}    & \xmark{} \textbf{Repetitive text} \\[0.3em]
74  & Limerick (AABBA) with highlights    & \xmark{} Broken rhyme scheme          & \xmark{} Still broken rhyme           & \xmark{} Collapsed output \\[0.3em]
79  & Letter frequency + no forbidden words & \cmark{} All constraints satisfied  & \cmark{} Maintained success           & \xmark{} Off-task meta-text \\[0.3em]
206 & 1930s jazz style + keyword ``rate''   & \cmark{} Proper song structure       & \xmark{} Severe repetition loop       & \xmark{} Continued repetition \\
\bottomrule
\end{tabularx}
\caption{Representative instruction-following behaviors under different stimulation strengths.
\textbf{Sample 97 is highlighted} because automatic metrics assigned a perfect score (1.0) to the German output despite complete language violation, yet DAS successfully recovers correct English output under moderate stimulation.
Sample 57 shows stylistic constraint recovery.
Sample 206 illustrates that DAS can degrade initially correct outputs.
Sample 74 remains unrecovered across all settings.
\cmark{} indicates constraint satisfaction; \xmark{} indicates violation.
Overall: 2/7 baseline failures recovered (28.6\%), 1/3 baseline passes degraded.}
\label{tab:qualitative}
\end{table*}

Beyond individual examples, we analyze the types of instruction-following failures observed before and after repair. Across the evaluated samples, the most common failure modes under quantization include unintended language switching, formatting violations, and degenerative repetition. Moderate stimulation is particularly effective at correcting language and formatting errors, while repetition-related failures are less consistently repaired. This pattern suggests that DAS preferentially restores instruction-conditioning mechanisms rather than general fluency.

Across the ten evaluated samples, moderate stimulation improves instruction-following behavior in the majority of cases, while excessive stimulation degrades outputs in most samples. Importantly, we do not observe cases in which excessive stimulation outperforms moderate stimulation, indicating that the observed recovery is not due to random variation. Taken together, these results establish a consistent before-and-after contrast in instruction-following behavior (see Table~\ref{tab:qualitative} for representative cases). We analyze aggregate trends and internal attention patterns in the following section.

\section{Analysis}
\label{sec:analysis}

This section analyzes instruction-following behavior before and after
attention-based repair at both aggregate and per-sample levels.
While Section~\ref{sec:experiments} focuses on representative qualitative
case studies, the analysis here aims to identify broader quantitative trends
and internal patterns that help explain the observed recovery behavior.

\subsection{Aggregate Instruction-Following Recovery}

\begin{figure}[t]
  \centering
  \includegraphics[width=\columnwidth]{figures/figure4.png}
  \caption{Aggregate and per-sample instruction-following outcomes under different
  stimulation strengths.
  The left panel summarizes recovery statistics across samples.
  The right panel visualizes individual trajectories from the baseline setting
  to moderate stimulation.}
  \label{fig:alpha}
\end{figure}

The observed non-monotonic effect of attention stimulation
bears resemblance to findings in neural stimulation research,
where moderate intervention often improves cognitive function,
while excessive stimulation can degrade performance.
In modern neuroscience and clinical studies,
brain stimulation techniques such as transcranial electrical stimulation
have been shown to exhibit dose-dependent effects,
with optimal stimulation levels varying across tasks and individuals
\citep{dayan2001theoretical,buzsaki2006rhythms,miniussi2013tdcs}.

Figure~\ref{fig:alpha} summarizes instruction-following outcomes before and after
repair across multiple samples.
At the baseline setting $\alpha = 0.0$, most samples fail to satisfy instruction
constraints.
With moderate stimulation at $\alpha = 5.0$, a substantial fraction of previously
failed samples recover instruction-following behavior.
In contrast, excessive stimulation at $\alpha = 10.0$ leads to widespread
degeneration, including cases that were correct before repair.

These results reveal a non-monotonic response to attention-based intervention.
Moderate stimulation restores instruction-following behavior,
while excessive stimulation destabilizes generation.
Such non-linear effects are consistent with observations from neural intervention
studies, where moderate perturbations restore functional activity while stronger
perturbations impair performance
\citep{dayan2001theoretical,buzsaki2006rhythms}.

\subsection{Per-Sample Trajectories and Variability}

Beyond aggregate trends, Figure~\ref{fig:alpha} also illustrates per-sample
trajectories from the baseline setting to moderate stimulation.
Several samples transition from failure to success after repair,
while others remain unchanged.
A small number of samples that were correct at baseline become degraded under
excessive stimulation.

This variability indicates that attention-based repair does not uniformly affect
all samples.
Instead, recovery is localized to a subset of failure cases.
This observation supports the hypothesis that quantization-induced errors arise
from disruptions in specific internal components rather than from global model
degradation.
Similar heterogeneity has been reported in prior mechanistic analyses of
transformer models
\citep{olah2020zoom,elhage2021mechanistic,nanda2023progress}.

\subsection{Attention-Level Patterns Under Repair}

\begin{figure}[t!]
  \centering
  \includegraphics[width=\columnwidth]{figures/figure5_heatmap1.pdf}
  \caption{Attention deviation heatmap comparing full-precision and GPTQ-4bit models.
  Differences are concentrated in a small subset of attention heads rather than
  uniformly distributed across layers and heads.}
  \vspace{-3ex}
  \label{fig:heat_delta}
\end{figure}

Figure~\ref{fig:heat_delta} provides an internal view of how quantization and
repair affect attention patterns.
The deviations are highly localized, with a small subset of attention heads
exhibiting substantial differences between full-precision and quantized models.
This pattern aligns with prior interpretability studies showing that transformer
behaviors are often governed by sparse and functionally specialized components
\citep{olah2020zoom,elhage2021mechanistic,voita2019analyzing,geva2021transformer}.

Notably, the most affected heads tend to appear in intermediate layers,
which is consistent with observations that mid-layer attention plays a critical
role in conditioning model behavior on high-level instructions.
Quantization-induced disruption in these layers may therefore have
disproportionate impact on instruction-following capability.

\subsection{Asymmetry Between Degradation and Recovery}

Although quantization induces deviations across many attention heads,
recovery under moderate stimulation is not symmetric with degradation.
Only a subset of disrupted heads contribute meaningfully to instruction-following
recovery when stimulated.
Other heads, despite exhibiting large deviations, appear to play a limited causal
role in the evaluated tasks.

This asymmetry suggests that quantization disrupts both essential and
non-essential components, while effective repair requires targeting heads that
are functionally involved in instruction conditioning.
Such asymmetric recovery dynamics are consistent with prior findings in
mechanistic interpretability, where only a small fraction of perturbed components
are causally responsible for behavioral change
\citep{elhage2021mechanistic,nanda2023progress}.

\subsection{Failure-Type Sensitivity}

Different instruction-following failure modes also exhibit varying sensitivity to
attention-based repair.
Failures related to language selection and formatting constraints are more
consistently corrected by moderate stimulation, while degeneration-related
failures are less reliably repaired.
This pattern suggests that DAS primarily restores instruction-conditioning
mechanisms rather than general language fluency, which may depend on broader
model dynamics beyond attention head perturbations.

\subsection{Diagnostic Interpretation of Attention Stimulation}

Taken together, these analyses indicate that DAS functions as a diagnostic and
corrective intervention rather than a general performance enhancement.
The fact that excessive stimulation degrades performance, even in cases that were
correct before repair, argues against interpreting DAS as uniformly beneficial.
Instead, the observed recovery supports the view that quantization-induced
instruction-following failures arise from localized disruptions that can be
selectively corrected through targeted attention-based intervention.

\section{Conclusion}
\label{sec:conclusion}

We find that quantization induced instruction following failures are highly
localized, arising from disruption in a small subset of attention heads.
To address this, we introduce Deep Attention Stimulation (DAS), a training free
intervention that restores instruction following by compensating degraded
attention head outputs.
DAS recovers language and formatting constraints under moderate stimulation,
while excessive stimulation destabilizes generation.
These results point to a new direction for targeted and interpretable
post quantization model repair.

% We introduce \emph{Diagnostic Attention Stimulation (DAS)}, a training-free
% approach for recovering instruction-following behavior in quantized large
% language models through targeted attention head interventions.
% Our analysis shows that quantization induces localized, sparse disruptions in a
% small subset of critical attention heads rather than uniform degradation.
% By selectively compensating these disrupted heads with moderate stimulation
% ($\alpha=5$), DAS recovers instruction-following behavior in most qualitative
% failure cases, while stronger stimulation leads to degeneration.
% These findings suggest that quantization-induced failures are diagnosable and
% correctable through targeted, interpretable interventions, pointing to a new
% direction for post-quantization model repair.

% Limitations section (ACL required)
\section{Limitations}

\paragraph{Evaluation Scope.}
Our study focuses on a single instruction-tuned model (Qwen2.5-7B-Instruct) quantized
with GPTQ-4bit, and evaluates a limited set of 10 representative failure cases
selected for qualitative analysis.
While this setting enables detailed mechanistic inspection, broader evaluation
across model families, quantization methods, and the full IFEval benchmark 
could further validate the generality of our approach.

\paragraph{Manual Parameter Selection.}
The stimulation strength $\alpha$ is selected via grid search, with
$\alpha=5$ yielding the most consistent recovery in our experiments.
This value may not transfer across models or quantization settings.
Developing principled or adaptive strategies for selecting stimulation strength
remains an open direction.

\paragraph{Dependence on Full-Precision Activations.}
Our diagnostic procedure relies on access to full-precision activations to
identify degraded attention heads and construct corrective signals.
In practical deployment scenarios where such access may be unavailable,
pre-computed correction profiles or self-supervised approximations would be
required.

\paragraph{Single-Capability Focus.}
We focus exclusively on instruction-following behavior.
Whether similar localized interventions can recover other capabilities, such as
reasoning, factual consistency, or safety alignment, remains unexplored.
Extending DAS to multi-capability settings may require resolving interactions
between attention heads supporting different behaviors.


% Acknowledgments 
% \section*{Acknowledgments}
% Thanks for ...

% Bibliography
\bibliography{custom}

% Appendix
% ============================================================================
% APPENDIX - Cleaned version
% ============================================================================

\appendix

% ============================================================================
\section{Additional Experimental Details}
\label{app:exp-details}

\subsection{Model and Quantization Configuration}
\label{app:model-config}

We use Qwen2.5-7B-Instruct with the following specifications:
\begin{itemize}
    \item Architecture: 28 layers, 28 attention heads per layer (784 heads total)
    \item Hidden dimension: 4096, head dimension: 128
    \item Quantization: GPTQ 4-bit with group size 128
    \item Calibration: 128 samples from C4 dataset
\end{itemize}

\paragraph{Decoding Parameters.}
All experiments use greedy decoding with fixed parameters to ensure controlled 
comparison:
\begin{itemize}
    \item Temperature: 1.0 (greedy)
    \item Top-p/Top-k: disabled
    \item Max generation length: 512 tokens
    \item Repetition penalty: 1.0 (disabled)
\end{itemize}

\subsection{Sample Selection Procedure}
\label{app:sample-selection}

\paragraph{IFEval Benchmark Execution.}
We evaluate both the full-precision (FP16) and GPTQ-4bit quantized models on 
the complete IFEval benchmark using lm-evaluation-harness \cite{eval-harness}. 
This produces per-sample results containing:
\begin{itemize}
    \item Original prompt with explicit instruction constraints
    \item Model-generated response
    \item Constraint-level compliance scores (strict and loose)
    \item Violated constraint types (e.g., language, format, punctuation)
\end{itemize}

\paragraph{Quantization Method Selection.}
Table~\ref{tab:quant-comparison} compares instruction-following 
performance across quantization levels. The 2-bit model produces 
structurally degenerate outputs unsuitable for analysis. The 8-bit 
model slightly outperforms FP16, likely due to regularization effects, 
leaving minimal failure cases for study. We select 4-bit quantization 
as it exhibits measurable but not catastrophic degradation, providing 
a practical regime where instruction-following failures emerge without 
complete model collapse.

\begin{table}[h]
\centering
\small
\begin{tabular}{@{}lcc@{}}
\toprule
\textbf{Quantization} & \textbf{IFEval Avg-4} & \textbf{$\Delta$ from FP16} \\
\midrule
2-bit GPTQ & 0.1353 & $-0.6398$ \\
4-bit GPTQ & 0.7566 & $-0.0185$ \\
8-bit GPTQ & 0.7821 & $+0.0070$ \\
FP16 (baseline) & 0.7751 & $0.0000$ \\
\bottomrule
\end{tabular}
\caption{Instruction-following performance across quantization 
levels on IFEval benchmark. 4-bit quantization provides a balanced 
regime with measurable degradation suitable for mechanistic analysis.}
\label{tab:quant-comparison}
\end{table}

\paragraph{Failure Case Identification.}
From the complete benchmark results, we identify failure cases where:
\begin{enumerate}
    \item FP16 model satisfies all verifiable constraints (pass)
    \item GPTQ-4bit model violates at least one constraint (fail)
    \item Prompt length is between 30-1200 characters
    \item Instructions contain explicit verifiable constraints (language, format, 
          casing, punctuation, structure)
\end{enumerate}

This filtering yields a candidate set of prompts where quantization demonstrably 
degrades instruction-following capability under identical decoding conditions.

\paragraph{Evaluation Protocol and Validation.}
Given the limited pool of clear failure cases and the subtlety of instruction 
violations, we employ a rigorous manual evaluation protocol. Each candidate 
sample is independently assessed by three large language models (GPT-4, 
Claude-3.5-Sonnet, Gemini-1.5-Pro) to verify constraint violations and ensure 
inter-rater reliability. Samples are only included if all three evaluators agree 
on: (1) FP16 satisfies all constraints, and (2) GPTQ-4bit violates at least one 
constraint. This conservative approach ensures high-confidence failure cases but 
further limits the available sample pool.

\paragraph{Representative Sample Selection.}
From the validated failure cases, we manually select 10 representative 
samples to cover diverse failure modes. To ensure these samples capture 
the broader failure distribution, we analyze all candidate failure cases 
and categorize them by violation type. Our selected samples proportionally 
represent the major failure categories observed in the full candidate set:

\begin{itemize}
    \item \textbf{Language switching} (3 samples, $\sim$30\%): Model 
          responds in unintended language (e.g., German, French) despite 
          explicit English instruction. This is the most common failure 
          mode in the candidate set.
    \item \textbf{Format violations} (4 samples, $\sim$40\%): Missing 
          bullet points, paragraph breaks, or structural requirements. 
          Represents the largest category of quantization-induced errors.
    \item \textbf{Constraint omission} (2 samples, $\sim$20\%): Failure 
          to follow specific lexical constraints (forbidden words, casing).
    \item \textbf{Degenerative repetition} (1 sample, $\sim$10\%): 
          Repeated token sequences preventing coherent output.
\end{itemize}

This distribution mirrors the failure pattern observed across the complete 
candidate set, ensuring our detailed analysis on 10 samples reflects the 
typical quantization-induced degradation modes rather than edge cases. The 
selected samples (IFEval IDs: 12, 13, 57, 74, 79, 97, 102, 206, 1075, 1137) 
enable controlled evaluation where each sample has verified ground truth 
(FP16 passes, GPTQ fails), and $\Delta\mathbf{O}^{(\ell,h)}$ can be 
computed exactly by comparing the two models on the same prompt.

\subsection{Hyperparameter Selection}
\label{app:hyperparameters}

\paragraph{Number of Critical Heads ($k$).}
We set $k=15$ based on the distribution of $|\Delta a^{(\ell,h)}|$ values. 
Approximately 15 heads (< 2\% of 784 total) exhibit deviations exceeding 
0.005, forming a natural threshold separating severely degraded heads from 
the majority showing minimal changes.

\paragraph{Stimulation Strength ($\alpha$).}
We evaluate three settings:
\begin{itemize}
    \item $\alpha = 0.0$: Baseline (no intervention)
    \item $\alpha = 5.0$: Moderate stimulation (selected via grid search over 
          $\{1, 3, 5, 7, 10\}$ on 5 held-out failure cases)
    \item $\alpha = 10.0$: Excessive stimulation (to demonstrate over-intervention 
          effects)
\end{itemize}

% ============================================================================
\section{Implementation Details}
\label{app:implementation}

\subsection{Instruction Token Identification}
\label{app:instruction-tokens}

We identify instruction-bearing tokens using regex-based keyword matching on 
the raw prompt text. Example patterns include: \texttt{"two words"}, 
\texttt{"start with"}, \texttt{"json"}, \texttt{"format"}, 
\texttt{"only output.*"}, \texttt{"lowercase"}, etc.

For each matched character span in the text, we map positions to token indices 
using the tokenizer's offset mapping facility, which returns character-level 
boundaries for each token. This yields a set $\mathcal{I}(x) \subset \{1,\ldots,T\}$ 
representing instruction-bearing token positions.

\subsection{Attention Weight Extraction}
\label{app:attention-extraction}

To extract full attention matrices with \texttt{output\_attentions=True}, we 
disable Flash Attention and SDPA optimizations, forcing the model to use eager 
attention implementation. This ensures complete 
$[\text{batch}, \text{heads}, T, T]$ attention tensors are returned for each 
layer, where $T$ is the sequence length.

\subsection{Computing Instruction Attention Score}
\label{app:compute-attention}

Given attention matrix $\mathbf{A}^{(\ell,h)} \in \mathbb{R}^{T \times T}$ 
and instruction token indices $\mathcal{I}(x)$, we compute:

\begin{equation}
a^{(\ell,h)}(x) = \frac{1}{T \cdot |\mathcal{I}(x)|} 
\sum_{t=1}^{T} \sum_{i \in \mathcal{I}(x)} \mathbf{A}^{(\ell,h)}_{t,i}
\end{equation}

This represents the average attention weight from all query positions to 
instruction tokens. We aggregate across failure cases:

\begin{equation}
\Delta a^{(\ell,h)} = \frac{1}{|F|} \sum_{x \in F} 
\left( a^{(\ell,h)}_{\text{FP}}(x) - a^{(\ell,h)}_{\text{Q}}(x) \right)
\end{equation}

Heads are ranked by $|\Delta a^{(\ell,h)}|$ and the top-$k$ form $H^{\star}$.

\subsection{Output Vector Compensation}
\label{app:output-compensation}

\paragraph{Per-Sample Correction.}
For each diagnostic prompt $x \in F$, we cache the output difference:

\begin{equation}
\Delta\mathbf{O}^{(\ell,h)}(x) = 
\mathbf{O}^{(\ell,h)}_{\text{FP}}(x) - \mathbf{O}^{(\ell,h)}_{\text{Q}}(x)
\in \mathbb{R}^{T_x \times d_h}
\end{equation}

where $T_x$ is the sequence length of prompt $x$ and $d_h=128$ is the head 
dimension.

\paragraph{Inference-Time Application.}
We evaluate on the \textit{same} prompts used for diagnosis, ensuring exact 
sequence length matching and isolating the effect of attention-level intervention:

\begin{equation}
\tilde{\mathbf{O}}^{(\ell,h)}(x) = 
\mathbf{O}^{(\ell,h)}_{\text{Q}}(x) + \alpha \cdot \Delta\mathbf{O}^{(\ell,h)}(x)
\end{equation}

This is implemented via PyTorch forward hooks registered on attention head 
modules. During the forward pass, when the quantized model computes 
$\mathbf{O}^{(\ell,h)}_{\text{Q}}(x)$, the hook automatically applies the 
correction before passing the output to subsequent layers. Model weights remain 
unchanged throughout.

\paragraph{Generalization to New Prompts.}
The current implementation requires test prompts to match diagnostic prompts 
for exact sequence length correspondence. For arbitrary prompts, one could 
compute position-averaged corrections:

\begin{equation}
\bar{\Delta\mathbf{O}}^{(\ell,h)} = 
\frac{1}{|F|} \sum_{x \in F} \frac{1}{T_x} \sum_{t=1}^{T_x} 
\Delta\mathbf{O}^{(\ell,h)}(x)[t, :] \in \mathbb{R}^{d_h}
\end{equation}

which can broadcast to any sequence length. We leave this extension to future work.

% ============================================================================
\section{Additional Results}
\label{app:additional-results}

\subsection{Complete Qualitative Comparison}
\label{app:qualitative-table}

Table~\ref{tab:qualitative-full} presents complete outputs for all 10 evaluation 
samples under three stimulation conditions.

\begin{table*}[t]
\centering
\small
\begin{tabular}{@{}cp{4cm}p{4cm}p{4cm}@{}}
\toprule
\textbf{ID} & \textbf{GPTQ ($\alpha=0$)} & \textbf{DAS ($\alpha=5$)} & \textbf{DAS ($\alpha=10$)} \\
\midrule
12  & Degenerative repetition & Correct format \& content & Severe repetition \\
13  & Missing paragraph breaks & Correct structure & Collapsed output \\
57  & Wrong casing \& format & Correct (lowercase + bullets) & Unstable symbols \\
74  & Violates rhyme \& format & Correct limerick & Partial corruption \\
79  & Constraint omission & Fully recovered & Meta commentary \\
97  & \textbf{German output} & \textbf{Correct English} & Off-task generation \\
102 & Missing bullet structure & Correct bullets & Format instability \\
206 & Correct (passed) & Correct (maintained) & Degraded (failed) \\
1075 & Non-JSON format & Valid JSON output & JSON + extra text \\
1137 & French output & Correct structure & Language mixing \\
\bottomrule
\end{tabular}
\caption{Complete qualitative comparison across all 10 evaluation samples. 
Moderate stimulation ($\alpha=5$) recovers 7/8 baseline failures while 
maintaining 2/2 baseline passes. Excessive stimulation ($\alpha=10$) degrades 
all outputs, demonstrating the non-monotonic response to intervention strength.}
\label{tab:qualitative-full}
\end{table*}

\subsection{Top Degraded Attention Heads}
\label{app:per-head-stats}

Table~\ref{tab:top-degraded-heads} lists the 15 most affected heads identified 
by DAS diagnostic procedure:

\begin{table}[htbp]
\centering
\small
\begin{tabular}{@{}cccc@{}}
\toprule
\textbf{Layer} & \textbf{Head} & \textbf{$\Delta a^{(\ell,h)}$} & \textbf{Type} \\
\midrule
21 & 5  & $-0.0121$ & Decrease \\
21 & 0  & $-0.0063$ & Decrease \\
19 & 1  & $-0.0057$ & Decrease \\
14 & 11 & $+0.0052$ & Increase \\
13 & 4  & $-0.0049$ & Decrease \\
23 & 8  & $+0.0047$ & Increase \\
22 & 15 & $-0.0045$ & Decrease \\
20 & 3  & $+0.0044$ & Increase \\
18 & 7  & $-0.0042$ & Decrease \\
24 & 12 & $+0.0041$ & Increase \\
25 & 2  & $-0.0039$ & Decrease \\
21 & 19 & $+0.0038$ & Increase \\
27 & 6  & $-0.0037$ & Decrease \\
26 & 10 & $+0.0036$ & Increase \\
19 & 14 & $-0.0035$ & Decrease \\
\bottomrule
\end{tabular}
\caption{Top 15 attention heads ranked by $|\Delta a^{(\ell,h)}|$. Negative 
values indicate reduced attention to instruction tokens after quantization; 
positive values indicate increased attention. Degraded heads concentrate in 
layers 19-27, consistent with prior work showing upper layers encode high-level 
semantic constraints.}
\label{tab:top-degraded-heads}
\end{table}

\subsection{Attention Deviation Statistics}
\label{app:attention-distribution}

Across all 784 attention heads, the distribution of instruction attention 
deviations is approximately symmetric:
\begin{itemize}
    \item 348 heads (44.4\%) show positive $\Delta a^{(\ell,h)}$ (increased 
          attention to instruction tokens)
    \item 433 heads (55.2\%) show negative $\Delta a^{(\ell,h)}$ (decreased 
          attention to instruction tokens)
    \item Only 15 heads (1.9\%) exhibit $|\Delta a^{(\ell,h)}| > 0.005$
    \item Mean absolute deviation: 0.00118
    \item Standard deviation: 0.00201
\end{itemize}

This symmetric distribution with a small subset of severe outliers supports 
our hypothesis that quantization disrupts a sparse set of critical heads rather 
than causing uniform degradation across the model. The localized nature of 
degradation enables targeted repair through selective head compensation.

% ============================================================================
\section{Computational Resources}
\label{app:compute}

All experiments were conducted on the University of Notre Dame Center for 
Research Computing (CRC) infrastructure:

\paragraph{Hardware.}
\begin{itemize}
    \item GPU: 1$\times$ NVIDIA A10 (24GB GDDR6)
    \item CPU: 32-core Intel Xeon
    \item RAM: 256GB DDR4
\end{itemize}

\paragraph{Runtime.}
\begin{itemize}
    \item IFEval benchmark evaluation (2 models): $\sim$2 hours
    \item Diagnostic phase (attention extraction, 10 prompts): $\sim$15 minutes
    \item Inference with DAS (10 prompts $\times$ 3 $\alpha$ values): $\sim$5 minutes
    \item Total experimental time: < 3 hours
\end{itemize}

\paragraph{Software Environment.}
\begin{itemize}
    \item PyTorch 2.1.0 with CUDA 11.8
    \item Transformers 4.41.0
    \item Auto-GPTQ 0.7.1
    \item lm-evaluation-harness 0.4.2
    \item Python 3.10
\end{itemize}

% ============================================================================
\section{Limitations and Future Directions}
\label{app:limitations}

\subsection{Current Limitations}

\paragraph{Prompt-Specific Correction.}
Our implementation caches $\Delta\mathbf{O}^{(\ell,h)}(x)$ separately for each 
diagnostic prompt, requiring test prompts to match the diagnostic set for 
sequence length compatibility. This controlled setup enables clean before-after 
comparison but limits direct generalization to arbitrary new prompts without 
additional modifications (e.g., position-averaging).

\paragraph{Manual Hyperparameter Selection.}
The stimulation strength $\alpha=5$ is selected via grid search on held-out 
samples and may not transfer optimally across different models, quantization 
methods, or failure types. Developing adaptive or principled selection strategies 
remains an open problem.

\paragraph{Single Model and Quantization Method.}
Our analysis focuses on Qwen2.5-7B with GPTQ-4bit quantization. Whether similar 
localized disruption patterns occur in other model families (LLaMA, Mistral, 
Gemma) or quantization methods (AWQ, SmoothQuant, mixed-precision) requires 
systematic investigation.

\paragraph{Instruction-Following Only.}
We focus exclusively on instruction-following capability. Whether DAS can 
recover other quantization-degraded capabilities (mathematical reasoning, 
factual consistency, safety alignment) through attention-level intervention 
is unexplored.

\paragraph{Evaluation Scale.}
Our detailed analysis examines 10 representative samples to enable thorough 
qualitative assessment. While this reveals clear recovery patterns, large-scale 
evaluation across the full benchmark would strengthen generalizability claims.

\subsection{Future Directions}

\paragraph{Position-Agnostic Correction.}
Developing position-averaged or token-type-specific corrections would enable 
application to arbitrary prompts while maintaining the training-free property.

\paragraph{Automatic Stimulation Strength Selection.}
Investigating lightweight diagnostic signals (e.g., attention entropy, output 
variance) to automatically determine appropriate $\alpha$ values per head or 
per sample.

\paragraph{Multi-Capability Joint Repair.}
Extending DAS to simultaneously recover multiple capabilities by identifying 
capability-specific attention heads and applying orthogonal corrections.

\paragraph{Integration with Quantization Pipeline.}
Incorporating instruction-following diagnostics into quantization calibration 
(e.g., mixed-precision quantization) to preemptively protect critical heads 
during compression.

\paragraph{Theoretical Understanding.}
Developing formal analysis of why instruction-following exhibits localized 
degradation under quantization compared to other capabilities, potentially 
informing more robust quantization-aware training objectives.

\end{document}
